\subsection{Aritmética de precisão arbitrária}

Para representar números muito grandes, podemos usar um vector de int em que cada elemento representa um dígito e isso na base $10^9$

\inccode{algtn/arbitraria-aritm.cpp}

Também podemos fatorar em fatores primos e trabalhar com eles.

Outra ideia é usar o algoritmo de Garner.

Uma ideia para variáveis de ponto flutuante é guardar um int com o expoente e um double com um valor em [0.1, 1)

\subsubsection{IO do int128}

\inccode{algtn/int128.cpp}
